\section*{Bab \ref{ch:b3:many-electron-atoms} --- Atom Berelektron Banyak dan Simbol Suku}

\begin{solution}{exo:b3:many-electron-atoms:1}
\[
\Psi = \frac{1}{\sqrt6}\begin{vmatrix}1s\alpha(1) & 1s\beta(1) & 2s\alpha(1)\\
1s\alpha(2) & 1s\beta(2) & 2s\alpha(2)\\ 1s\alpha(3) & 1s\beta(3) & 2s\alpha(3)
\end{vmatrix}.
\]
$1s$ hanya menyediakan dua \omterm{def:b3:many-electron-atoms:spin-orbital}{spin-orbital}, $1s\alpha$ dan $1s\beta$; elektron
ketiga akan mengulang salah satunya sehingga dua kolom menjadi sama:
determinannya nol.
\end{solution}

\begin{solution}{exo:b3:many-electron-atoms:2}
$\binom62 = 15$, $\binom63 = 20$, $\binom{10}{2} = 45$. $p^3$: ${}^4$S (4) +
${}^2$D (10) + ${}^2$P (6) = 20. $d^2$: ${}^3$F (21) + ${}^3$P (9) + ${}^1$G (9) +
${}^1$D (5) + ${}^1$S (1) = 45.
\end{solution}

\begin{solution}{exo:b3:many-electron-atoms:3}
N ($2p^3$, setengah penuh): \termsym{4}{S}{3/2}. O ($2p^4$, lebih dari
setengah): \termsym{3}{P}{2}. F ($2p^5$): \termsym{2}{P}{3/2}. \ce{Ti^2+} ($3d^2$):
\termsym{3}{F}{2}. \ce{Fe^3+} ($3d^5$, setengah penuh, semua spin sejajar,
$L = 0$): \termsym{6}{S}{5/2}.
\end{solution}

\begin{solution}{exo:b3:many-electron-atoms:4}
${}^2$D: $J = \frac52$ (6 keadaan) dan $\frac32$ (4); $6 + 4 = 10 = 5 \times 2$.
${}^3$P: $J = 2$ (5), 1 (3), 0 (1); $5 + 3 + 1 = 9 = 3 \times 3$.
\end{solution}

\begin{solution}{exo:b3:many-electron-atoms:5}
Interval $16.4$ ($J = 1 - 0$) dan $27.0$ ($2 - 1$): rasionya 1.65, bukan 2.
$A = 16.4/1 = \qty{16.4}{cm^{-1}}$ dari interval pertama, $27.0/2 = \qty{13.5}{cm^{-1}}$
dari interval kedua. Satu nilai $A$ tidak cocok untuk keduanya: tingkat-tingkat
itu sedikit bercampur dengan tingkat-tingkat suku lain (${}^1$D, ${}^1$S), dan
\omterm{def:b3:many-electron-atoms:russell-saunders}{kopling Russell--Saunders} hanyalah sebuah pendekatan.
\end{solution}

\begin{solution}{exo:b3:many-electron-atoms:6}
$\lambda = 10^7/\tilde\nu$: \qty{589.76}{nm} (D$_1$) dan \qty{589.16}{nm}
(D$_2$). Jaraknya $\qty{17.20}{cm^{-1}} = \qty{2.13}{meV}$. $E(\frac32) -
E(\frac12) = \frac32A$, sehingga $A = \qty{11.47}{cm^{-1}}$.
\end{solution}

\begin{solution}{exo:b3:many-electron-atoms:7}
$3p \to 3s$: diizinkan ($\Delta l = -1$). $3d \to 3s$: terlarang ($\Delta l = -2$).
$3d \to 3p$: diizinkan. $4s \to 3s$: terlarang ($\Delta l = 0$, tanpa perubahan
paritas). $4s \to 3p$: diizinkan.
\end{solution}

\begin{solution}{exo:b3:many-electron-atoms:8}
Rata-rata ${}^3$P: $(5 \times 169086.76 + 3 \times 169086.84 + 169087.83)/9 =
\qty{169086.91}{cm^{-1}}$. Selisihnya terhadap ${}^1$P: \qty{2047.98}{cm^{-1}} $=
\qty{0.254}{eV}$, sehingga $K(1s,2p) = \qty{0.127}{eV}$, tiga kali lebih kecil
daripada $K(1s,2s)$. \omterm{def:b3:many-electron-atoms:exchange}{Integral pertukaran} mengukur tumpang-tindih kerapatan
kedua orbital; orbital $2s$ menembus daerah $1s$ (orbital itu mempunyai
kerapatan di dekat inti), sedangkan orbital $2p$ lenyap di inti dan
tumpang-tindihnya lebih kecil.
\end{solution}

\begin{solution}{exo:b3:many-electron-atoms:9}
$M_L = 4$ yang terbesar memerlukan kedua elektron dalam $m_l = 2$, dengan spin
berlawanan: sebuah singlet, ${}^1$G (9 keadaan). Berikutnya, $M_L = 3$, berasal
dari $m_l = 2, 1$, dengan $M_S = 1$ dimungkinkan: ${}^3$F (21). Sisa $M_L = 2$
hanya dengan $M_S = 0$: ${}^1$D (5). Sisa $M_L = 1$ dengan $M_S = 1$: ${}^3$P (9).
Tinggal satu keadaan dengan $M_L = M_S = 0$: ${}^1$S. Jumlahnya 45.
\end{solution}

\begin{solution}{exo:b3:many-electron-atoms:10}
Pada $\mathbf r_1 = \mathbf r_2 = \mathbf r$ fungsi triplet bernilai
$\frac{1}{\sqrt2}[1s(\mathbf r)2s(\mathbf r) - 2s(\mathbf r)1s(\mathbf r)] = 0$;
fungsi singlet bernilai $\sqrt2\,1s(\mathbf r)2s(\mathbf r) \ne 0$. Dalam
triplet, kedua elektron tidak pernah ditemukan di titik yang sama dan
rata-rata lebih berjauhan, sehingga tolakannya lebih kecil:
$E_+ - E_- = 2K > 0$.
\end{solution}

\begin{solution}{exo:b3:many-electron-atoms:11}
$d^8$ lebih dari setengah penuh: koplingnya terbalik, dengan $J = L + S = 4$
terendah. Interval: $E(3) - E(4) = 1360.7$ dan $E(2) - E(3) = 908.9$; rasionya
1.50, dibandingkan dengan $4/3 = 1.33$ menurut Landé. $|A| = 1360.7/4 = \qty{340}{cm^{-1}}$
($A < 0$). \omterm{def:b3:many-electron-atoms:spin-orbit}{Kopling spin--orbit} jauh lebih besar dalam ion yang lebih berat ini
daripada dalam karbon (\qty{16}{cm^{-1}}).
\end{solution}

\begin{solution}{exo:b3:many-electron-atoms:12}
$\sum_J(2J+1)[J(J+1) - L(L+1) - S(S+1)] = 0$ (keadaan-keadaan yang berjumlah $(2L+1)(2S+1)$
itu dapat dihitung dalam basis mana pun, dan $\sum M_LM_S = 0$ atas
keadaan-keadaan tak berkopling). Untuk ${}^3$P: $1(0 - 4) + 3(2 - 4) + 5(6 - 4) = -4 - 6 + 10 = 0$.
Rata-rata berbobot karbon adalah $(0 + 3 \times 16.42 + 5 \times 43.41)/9 = \qty{29.6}{cm^{-1}}$:
energi yang akan dimiliki suku itu tanpa \omterm{def:b3:many-electron-atoms:spin-orbit}{kopling spin--orbit}.
\end{solution}

\begin{solution}{pb:b3:many-electron-atoms:1}
\textbf{1.} $1s\alpha$, $1s\beta$, $2s\alpha$, $2s\beta$.
\textbf{2.} $|1s\alpha\,2s\alpha|$, $|1s\alpha\,2s\beta|$, $|1s\beta\,2s\alpha|$,
$|1s\beta\,2s\beta|$.
\textbf{3.} $|1s\alpha\,2s\alpha| = \frac{1}{\sqrt2}[1s(1)2s(2) - 2s(1)1s(2)]\,
\alpha(1)\alpha(2)$, dengan menjabarkan determinan $2\times2$ itu.
\textbf{4.} Jumlahnya adalah $\frac{1}{\sqrt2}[1s(1)2s(2) - 2s(1)1s(2)]\cdot
\frac{1}{\sqrt2}[\alpha(1)\beta(2) + \beta(1)\alpha(2)]$ (dikalikan $\sqrt2$, lalu
dinormalisasi); selisihnya memberikan $\frac{1}{\sqrt2}[1s(1)2s(2) +
2s(1)1s(2)]\cdot\frac{1}{\sqrt2}[\alpha(1)\beta(2) - \beta(1)\alpha(2)]$.
\textbf{5.} Fungsi ruang simetris dengan fungsi spin antisimetris ($S = 0$);
fungsi ruang antisimetris dengan ketiga fungsi spin simetris ($S = 1$).
\textbf{6.} $2S + 1 = 1$ dan 3: satu dan tiga keadaan dengan energi yang sama.
\textbf{7.} $[\ce{Ar}]\,3d^2$; $\binom{10}{2} = 45$ determinan.
\textbf{8.} Maksimum $S = 1$, maksimum $L = 2 + 1 = 3$, kurang dari
setengah penuh: \termsym{3}{F}{2}, tingkat pada 0 dalam data.
\textbf{9.} $21 + 9 + 9 + 5 + 1 = 45$.
\textbf{10.} Normal ($J = 2$ terendah). Interval 184.9 dan 235.5: rasionya
1.27 dibandingkan dengan $4/3 = 1.33$; aturan itu berlaku dengan selisih
kurang dari 5\,\%.
\textbf{11.} ${}^3$F: $(5 \times 0 + 7 \times 184.9 + 9 \times 420.4)/21 =
\qty{241.8}{cm^{-1}}$; ${}^3$P: $(10538.4 + 3 \times 10603.6 + 5 \times
10721.2)/9 = \qty{10661.7}{cm^{-1}}$.
\textbf{12.} $15B = \qty{10419.9}{cm^{-1}}$, $B = \qty{695}{cm^{-1}}$.
\textbf{13.} $\Delta S = 0$.
\textbf{14.} $\Delta S = 1$ terlarang, dan $1s2s \to 1s^2$ mempunyai $\Delta l = 0$
(tanpa perubahan paritas; ${}^3$S$_1 \to {}^1$S$_0$ juga mempunyai $\Delta L = 0$
antara dua keadaan S). Keadaan itu metastabil: umurnya jauh lebih panjang
daripada keadaan tereksitasi biasa.
\textbf{15.} $169086.91 - 159855.97 = \qty{9230.94}{cm^{-1}}$: \qty{1083.3}{nm},
di inframerah dekat.
\textbf{16.} $171134.89 - 166277.44 = \qty{4857.45}{cm^{-1}}$: \qty{2058.7}{nm}.
\textbf{17.} $\qty{171134.89}{cm^{-1}}$: \qty{58.43}{nm}, di ultraviolet vakum
(diserap oleh udara).
\textbf{18.} Setiap keluarga mempunyai garis-garisnya sendiri dan keadaan
terendahnya sendiri, dan keduanya tidak pernah terhubung oleh cahaya:
keduanya berperilaku seperti dua gas dengan dua spektrum.
\textbf{19.} $E_\pm = E_{1s} + E_{2s} + J \pm K$ (singlet $+$, triplet $-$).
\textbf{20.} Triplet: fungsi ruangnya lenyap ketika kedua elektron bertemu,
sehingga tolakannya lebih lemah.
\textbf{21.} \qty{6421.47}{cm^{-1}}, $6421.47/8065.54 = \qty{0.796}{eV}$.
\textbf{22.} Selisihnya $\qty{2047.98}{cm^{-1}} = \qty{0.254}{eV}$, sehingga
$K(1s,2p) =
\qty{0.127}{eV}$.
\textbf{23.} $2s$ menembus daerah $1s$ dan bertumpang-tindih dengannya lebih
besar daripada $2p$.
\textbf{24.} Ya: dalam kedua konfigurasi, triplet ($S$ lebih besar) berada
paling rendah.
\textbf{25.} $K = \frac12 \times \qty{0.796}{eV}$: \textbf{\omterm{def:b3:many-electron-atoms:exchange}{integral pertukaran}
$K(1s,2s)$ helium adalah \qty{0.398}{eV}}, yang diukur sebagai setengah
selisih singlet--triplet.
\end{solution}
